\specialsection{Введение}

Спортивное программирование в наше время приобретает всё большую популярность.
Проводятся всероссийские и международные соревнования.
В них участники решают алгоритмические задачи, которые требуют от них не только теоретических знаний,
но и практических навыков программирования, среди которых важную роль играет умение быстро и эффективно писать код.

Общепринятым форматом проведения соревнований является контест, индивидуальный или командный,
в котором может быть любое количество участников.
Зачастую контест содержит набор задач разной сложности, на решение которых отводится вплоть до пяти часов времени.
Участники соревнуются в количестве решённых задач либо в суммарном количестве баллов, набранных за решения задач.

Стандартная задача по спортивному программированию обычно состоит из условия, описания формата входных и выходных данных,
ограничений по времени и памяти, а также набора тестов.
Для решения задачи требуется написать код на одном из разрешённых языков программирования,
который будет обрабатывать входные данные и выдавать результат в соответствии с условием задачи.
В условии, как правило, приводятся открытые примеры тестов с правильными ответами,
а полная и окончательная проверка решения выполняется на скрытых тестах, недоступных участнику во время соревнования.

Формат дуэли представляет собой соревнование между двумя участниками,
которые одновременно решают одну и ту же задачу и стараются сделать это быстрее друг друга.
Такой вид соревнования вызывает больший интерес и повышает вовлечённость участников,
так как они соревнуются не с абстрактным списком соперников, а с одним конкретным человеком.
Дуэль обычно представляет собой короткое соревнование с более простыми задачами,
что требует высокой концентрации и быстроты действий в течение меньшего промежутка времени.
Участники отрабатывают и проявляют навыки эффективного написания кода,
а также умение находить и исправлять ошибки в условиях стресса и ограниченного времени, чтобы решить задачу раньше соперника.

Существующие платформы для спортивного программирования в основном ориентированы на проведение контестов
и не предоставляют достаточной гибкости и возможности для организации дуэлей по пользовательским правилам.
Дополнительный интерес вызывает анализ поведения участников в дуэлях и выявление закономерностей,
которые могут быть использованы для выявления нечестных пользователей,
использующих сторонние инструменты для получения преимущества.
Особенно актуальной становится задача отслеживания несанкционированного применения инструментов искусственного интеллекта,
способных генерировать готовые фрагменты решений или полностью решать задачу за участника.
Оперативное выявление таких случаев необходимо для сохранения духа честного соревнования,
в котором результат определяется навыками, подготовкой и скоростью самого участника.

В связи с этим возникает необходимость в разработке платформы для проведения дуэлей по спортивному программированию,
поддерживающей различные форматы дуэлей и турниров, автоматическое тестирование решений участников,
а также сбор и анализ информации о поведении участников для выявления возможных нарушений правил.
